\usepackage[utf8]{inputenc}
\usepackage[T1]{fontenc}
\usepackage[a4paper,margin=25mm,headsep=5mm,headheight=12pt]{geometry}
\frenchspacing
\usepackage{amsmath,amssymb,amsfonts,amsthm}  
\usepackage{hyperref} 
\hypersetup{
    colorlinks, 
    citecolor=black,
    filecolor=black, 
    linkcolor=black,
    urlcolor=black
}
\usepackage{setspace}
\usepackage{graphicx}
\usepackage{subcaption}
\usepackage[separate-uncertainty=true]{siunitx} 
\sisetup{
  range-phrase=--,
  detect-all,
  output-decimal-marker={.},
  range-units=single,
  per-mode=reciprocal, 
}
\usepackage{url}
\usepackage{verbatim}
\usepackage[toc,acronym]{glossaries}
\usepackage{booktabs} 
\usepackage{fancyhdr} 
\usepackage{lipsum}
\usepackage{lastpage} 
\usepackage{float} 
\usepackage{listings} 
\usepackage{comment} 
\usepackage{tcolorbox}
\usepackage[ddmmyyyy]{datetime}
\usepackage[english]{babel} 
\usepackage{enumitem}

\newdateformat{usvardate}{\monthname[\THEMONTH] \THEDAY, \THEYEAR}


\captionsetup{width=0.9\textwidth,font=footnotesize,labelfont=bf,labelsep=colon}
\renewcommand{\dateseparator}{--}

\makeatletter

\setlength{\parindent}{0pt}
\setlength{\parskip}{2em}


\bibliographystyle{unsrt}

% MY OWN PACKAGES


\newcommand{\abs}[1]{\left| #1 \right|}

\newcommand{\expectation}[1]{\left< #1 \right>}

\newcommand{\scalar}[2]{\left< #1 | #2 \right> }

\newcommand{\derivate}[2]{\frac{\partial #1}{\partial#2}}

\newcommand{\expscalar}[2]{\left< #1 \left| #2\right| #1 \right>}

\newcommand{\optscalar}[3]{\left< #1 \left| #2\right| #3 \right>}

\usepackage{tikz}

\newcommand{\set}[1]{\left\{#1 \right\}}

\newcommand{\norm}[1]{\left\|\Vec{#1} \right\|}

\newcommand{\expect}[1]{\left\langle #1 \right\rangle}

\newcommand{\nuc}[2]{\text{${}^{#1}$#2}}

\newcommand{\uvec}[1]{\mathbf{\hat{#1}}}

\newcommand{\degr}[0]{\text{$\text{}^\circ$}}


\DeclareSIUnit\angstrom{\protect \text {Å}}

\DeclareSIUnit\atomicmassunit{u}

\newenvironment{question}
  {\begin{tcolorbox}[colframe=white!50!black, colback=white!10!white, title=\section*{Question}] }
  {\end{tcolorbox}}           

\newenvironment{answer}
  {\begin{tcolorbox}[colframe=white!50!black, colback=white!10!white, title=Answer] }
  {\end{tcolorbox}}     
  
\usepackage{tikz}
\usetikzlibrary{matrix, backgrounds}
\usepackage{amsmath}
\usepackage{xcolor}
\usepackage{empheq}

% whitening factor (0 = original, 0.5 = softer, 0.7 = very soft)
\def\alpha{0.7}

\pgfmathdeclarefunction{R}{1}{%
  \pgfmathparse{(1-\alpha)*min(200, abs(#1) * 255) + \alpha*255}%
}
\pgfmathdeclarefunction{G}{1}{%
  \pgfmathparse{(1-\alpha)*min(255, max(0, 4 * #1 * abs(1 - #1) * 255)) + \alpha*255}%
}
\pgfmathdeclarefunction{B}{1}{%
  \pgfmathparse{(1-\alpha)*((abs(1 - #1)) * 255) + \alpha*255}%
}

\newcommand{\getcolor}[1]{%
  \pgfmathsetmacro{\red}{R(#1)}%
  \pgfmathsetmacro{\green}{G(#1)}%
  \pgfmathsetmacro{\blue}{B(#1)}%
  \definecolor{currentcolor}{RGB}{\red, \green, \blue}%
}

